\section{Introduction}
Browser fingerprints summarize properties exposed by a client, including its user agent, language, timezone, display, and graphics environment. Web services use such observations to recognize clients or support authentication decisions, while the same mechanisms create privacy risks~\cite{laperdrix2020browser,alaca2016device,spooren2015harmful}. Their usefulness depends partly on the integrity of the reports: an attacker who controls a reporting surface may change what a service sees without changing the underlying device. Fingerprint spoofing and inconsistent evasive fingerprints are established problems~\cite{liu2022gummy,vastel2018fpscanner,venugopalan2025fpinconsistent}.

Android hybrid applications make this problem distinctive. Web content executes inside a host application that can observe the operating system and configure a WebView. A JavaScript report is therefore one of several partially overlapping accounts of the same environment. An independent system browser can provide another account. This creates an opportunity to ask whether reports are mutually compatible, but also a risk of treating legitimate differences between APIs and containers as manipulation. The Web-security question we study is: \emph{which cross-layer references help detect manipulated mobile Web reports, and what normal-alarm and missing-evidence costs do they introduce?}

A single-field rule illustrates the difficulty. A large reported memory value may exceed a particular device's plausible capacity while remaining below a global threshold. Conversely, a changed timezone may reflect an ordinary system-setting change rather than a manipulated webpage. A screen-height threshold may flag both an effective modification and a harmless layout expansion. The relevant evidence is often a relationship to a contextual reference, not whether an isolated value is unusual. Reference observations are themselves untrusted measurements: consistency cannot attest that all surfaces are truthful.

We develop \hg to study this tradeoff. Its App collector records Android Native, WebView-host, and in-app Web observations under a fixed contract. Checks retain unknown and execution-failure states rather than silently assigning benign values. Offline selection combines admitted conditions into a compact detector subject to normal-alarm and definite-output constraints; online evaluation uses only the current observations required by the frozen detector. An optional pairing path links the App to an independent browser. We distinguish the full App detector, individual mechanism checks, and the retained App-plus-browser extension throughout.

Our evaluation has four parts. First, a controlled benchmark compares the full App detector with four retrained ablations and a small decision tree. Second, historical normal-device records expose the cost of optimizing only controlled detections. Third, memory, timezone, and geometry studies test concrete reference mechanisms against simpler conditions. Fourth, paired observations locate the incremental benefit of a cross-endpoint timezone check. The datasets and model configurations have development-use history; the study is a reproducible characterization of these materials, not a new blind test.

The contributions are:
\begin{itemize}[leftmargin=*,itemsep=1pt,topsep=2pt]
\item An explicit three-view cross-verification framework for Android Native, the WebView host, and in-app Web observations, with state-aware checks that separate contextual references, unavailable evidence, and manipulation alerts.
\item A controlled component analysis: \full flags 105/126 modifications with 0/252 controlled-normal alarms; disabling the historical-normal alarm cap adds three detections but produces alarms on all historical normal records.
\item Mechanism and extension evidence that identifies both benefits and limits: device references recover missed memory modifications and tolerate tested normal changes, while the retained browser extension adds five timezone detections but leaves language modifications unresolved.
\end{itemize}
The findings support evaluating reference checks jointly with normal behavior and evidence availability. They do not establish a universally trusted observation layer, complete manipulation coverage, or a calibrated authentication policy.
